\chapter{Maximum Principles}

\section{Strong Maximum Principle}\label{hanlin:maxprin-section-2-1}

Throughout this section, $\Omega\subset\mathbb{R}^n$ is bounded and connected, and
\[
Lu \equiv a_{ij}(x)D_{ij}u + b_i(x)D_iu + c(x)u,
\]
where the coefficients are continuous on $\overline\Omega$ and, for some
$\lambda>0$,
\[
a_{ij}(x)\xi_i\xi_j \ge \lambda |\xi|^2
\qquad (x\in\Omega,\ \xi\in\mathbb{R}^n).
\]

\begin{lemma}\label{hanlin:maxprin-lemma-2-1}\dcref{Lemma 2.1}\group{ch02-maximum-principles}
\source{han-lin-lecture-notes:page-0029}
\lean{HanLinLectureNotes.Ch02.positive_operator_nonnegative_maximum_not_mem}
\leanok
Let $u\in C^2(\Omega)\cap C(\overline\Omega)$, with $Lu>0$ and $c\le0$ in
$\Omega$. A nonnegative maximum of $u$ on $\overline\Omega$ cannot occur in
$\Omega$.
\end{lemma}
\begin{proof}
\sketch At an interior nonnegative maximizer $x_0$, one has $Du(x_0)=0$ and
$D^2u(x_0)\le0$. Positivity of $(a_{ij}(x_0))$ and $c(x_0)u(x_0)\le0$ give
$Lu(x_0)\le0$, contradicting $Lu>0$.
\end{proof}

\begin{remark}\label{hanlin:maxprin-remark-2-2}\dcref{Remark 2.2}\group{ch02-maximum-principles}
\source{han-lin-lecture-notes:page-0029}
\lean{HanLinLectureNotes.Ch02.positive_operator_maximum_not_mem_of_zeroOrder_eq_zero}
\leanok
When $c\equiv0$, the sign restriction on the maximum is unnecessary; the same
observation applies to the analogous results below.
\end{remark}

\begin{lemma}[Weak maximum principle]\label{hanlin:maxprin-lemma-2-3}\dcref{Lemma 2.3}\group{ch02-maximum-principles}
\source{han-lin-lecture-notes:page-0029}
\lean{HanLinLectureNotes.Ch02.weak_maximum_principle}
\leanok
Let $u\in C^2(\Omega)\cap C(\overline\Omega)$ satisfy $Lu\ge0$ and $c\le0$ in
$\Omega$. If its maximum on $\overline\Omega$ is nonnegative, that maximum is
attained on $\partial\Omega$.
\end{lemma}
\begin{proof}
\sketch For $\epsilon>0$, put $w=u+\epsilon e^{\alpha x_1}$. Then
\[
Lw = Lu + \epsilon e^{\alpha x_1}(a_{11}\alpha^2 + b_1\alpha + c).
\]
Boundedness of $b_1,c$ and $a_{11}\ge\lambda$ permit a fixed large $\alpha$ with
\[
a_{11}\alpha^2+b_1\alpha+c>0
\]
throughout $\Omega$. Lemma~\ref{hanlin:maxprin-lemma-2-1} therefore yields
\[
\sup_{\Omega} u \le \sup_{\Omega} w \le \sup_{\partial \Omega} w^+
\le \sup_{\partial \Omega} u^+ + \epsilon \sup_{x \in \partial \Omega} e^{\alpha x_1}.
\]
Letting $\epsilon\downarrow0$ proves the assertion.
\end{proof}

If $c\le0$, the weak maximum principle gives uniqueness in
$C^2(\Omega)\cap C(\overline\Omega)$ for
\[
Lu=f\quad\text{in }\Omega,
\qquad
u=\varphi\quad\text{on }\partial\Omega,
\]
where $f\in C(\Omega)$ and $\varphi\in C(\partial\Omega)$.

\begin{remark}\label{hanlin:maxprin-remark-2-4}\dcref{Remark 2.4}\group{ch02-maximum-principles}
\source{han-lin-lecture-notes:page-0030}
\lean{HanLinLectureNotes.dirichlet_uniqueness_counterexamples}
\leanok
The boundedness of $\Omega$ and the hypothesis $c\le0$ are essential in general;
Dirichlet uniqueness can fail on an unbounded domain or when $c$ has positive
part.
\end{remark}

\begin{example}\label{hanlin:maxprin-example-2-5}\dcref{Example 2.5}\group{ch02-maximum-principles}
\source{han-lin-lecture-notes:page-0030}
\lean{HanLinLectureNotes.sineProduct_square_counterexample}
\leanok
For $\Omega=(0,\pi)^2$, the nonzero function $u(x,y)=\sin x\sin y$ satisfies
\[
\Delta u+2u=0\quad\text{in }\Omega,
\qquad
u=0\quad\text{on }\partial\Omega.
\]
\end{example}

\begin{theorem}[Hopf lemma]\label{hanlin:maxprin-theorem-2-6}\dcref{Theorem 2.6}\group{ch02-maximum-principles}
\source{han-lin-lecture-notes:page-0030,han-lin-lecture-notes:page-0031}
\lean{HanLinLectureNotes.Ch02.hopf_boundary_point_liminf_pos}
\leanok
Let $B\subset\mathbb{R}^n$ be an open ball, $x_0\in\partial B$, and
$u\in C^2(B)\cap C(B\cup\{x_0\})$. Suppose $Lu\ge0$, $c\le0$ in $B$, and
\[
u(x)<u(x_0)\quad(x\in B),
\qquad u(x_0)\ge0.
\]
For every outward direction $v$ satisfying $v\cdot n(x_0)>0$,
\[
\liminf_{t\to0^+}\frac{u(x_0)-u(x_0-tv)}{t}>0.
\]
\end{theorem}
\begin{proof}
\sketch On an interior ball tangent at $x_0$, use
$h(x)=e^{-\alpha|x|^2}-e^{-\alpha r^2}$. Uniform ellipticity gives $Lh>0$
for large $\alpha$. Lemma~\ref{hanlin:maxprin-lemma-2-1} applied to
$u+\epsilon h$, for a fixed sufficiently small $\epsilon>0$, forces its maximum
to $x_0$. The outward difference quotient then gives the stated strict
inequality.
\end{proof}
\begin{remark}\label{hanlin:maxprin-remark-2-7}\dcref{Remark 2.7}\group{ch02-maximum-principles}
\source{han-lin-lecture-notes:page-0030,han-lin-lecture-notes:page-0031}
\lean{HanLinLectureNotes.Ch02.hopf_boundary_outward_normal_fderiv_pos}
\leanok
If $u\in C^1(B\cup\{x_0\})$, then
\[
\frac{\partial u}{\partial \nu}(x_0) > 0.
\]
\end{remark}

\begin{theorem}[Strong maximum principle]\label{hanlin:maxprin-theorem-2-8}\dcref{Theorem 2.8}\group{ch02-maximum-principles}
\source{han-lin-lecture-notes:page-0031}
\lean{HanLinLectureNotes.Ch02.strong_maximum_principle}
\leanok
Let $u\in C^2(\Omega)\cap C(\overline\Omega)$ satisfy $Lu\ge0$ and $c\le0$ in
$\Omega$. Unless $u$ is constant, every nonnegative maximum on
$\overline\Omega$ occurs only on $\partial\Omega$.
\end{theorem}
\begin{proof}
\sketch Let $M\ge0$ be the maximum and
$\Sigma=\{x\in\Omega:u(x)=M\}$. If $\Sigma$ is nonempty but not all of
$\Omega$, connectedness produces a ball $B\Subset\Omega\setminus\Sigma$ tangent
to $\Sigma$ at some $x_0$. Thus
\[
u(x)<u(x_0)=M\quad(x\in B).
\]
Theorem~\ref{hanlin:maxprin-theorem-2-6} gives a positive outward derivative at
$x_0$, whereas the interior maximum gives $Du(x_0)=0$.
\end{proof}

\begin{corollary}[Comparison principle]\label{hanlin:maxprin-cor-2-9}\dcref{Corollary 2.9}\group{ch02-maximum-principles}
\source{han-lin-lecture-notes:page-0032}
Suppose $u\in C^2(\Omega)\cap C(\overline\Omega)$, $Lu\ge0$, $c\le0$ in
$\Omega$, and $u\le0$ on $\partial\Omega$. Then $u\le0$ in $\Omega$; moreover,
either $u<0$ there or $u\equiv0$.
\end{corollary}

\begin{corollary}\label{hanlin:maxprin-cor-2-10}\dcref{Corollary 2.10}\group{ch02-maximum-principles}
\source{han-lin-lecture-notes:page-0032}
Assume that $\Omega$ has the interior sphere property and
$u\in C^2(\Omega)\cap C^1(\overline\Omega)$ satisfies $Lu\ge0$ and $c\le0$ in
$\Omega$. If a nonconstant $u$ attains a nonnegative maximum at $x_0$, then
$x_0\in\partial\Omega$ and every outward direction $v$ at $x_0$ satisfies
\[
\frac{\partial u}{\partial v}(x_0)>0,
\]
\end{corollary}

\begin{corollary}\label{hanlin:maxprin-cor-2-11}\dcref{Corollary 2.11}\group{ch02-maximum-principles}
\source{han-lin-lecture-notes:page-0032}
Let $\Omega\subset\mathbb{R}^n$ be bounded with the interior sphere property,
and let $u\in C^2(\Omega)\cap C^1(\overline\Omega)$ solve
\[
Lu=f\quad\text{in }\Omega,
\qquad
\frac{\partial u}{\partial n}+\alpha u=\varphi
\quad\text{on }\partial\Omega,
\]
where $f\in C(\overline\Omega)$, $\varphi\in C(\partial\Omega)$, $c\le0$ in
$\Omega$, and $\alpha\ge0$ on $\partial\Omega$. The solution is unique if
$c\not\equiv0$ or $\alpha\not\equiv0$; if both vanish identically, it is unique
up to an additive constant.
\end{corollary}
\begin{proof}
\sketch It suffices to treat a homogeneous solution:
\[
Lu=0\quad\text{in }\Omega,
\qquad
\frac{\partial u}{\partial n}+\alpha u=0
\quad\text{on }\partial\Omega.
\]
If $c\not\equiv0$ or $\alpha\not\equiv0$, a nonzero constant is impossible; for
a nonconstant solution, Corollary~\ref{hanlin:maxprin-cor-2-10} contradicts the
boundary condition at a positive maximum. Applying the same argument to $-u$
gives $u\equiv0$. If $c\equiv\alpha\equiv0$, Theorem~\ref{hanlin:maxprin-theorem-2-8}
and Corollary~\ref{hanlin:maxprin-cor-2-10} instead force every homogeneous
solution to be constant.
\end{proof}

\begin{theorem}\label{hanlin:maxprin-theorem-2-12}\dcref{Theorem 2.12}\group{ch02-maximum-principles}
\source{han-lin-lecture-notes:page-0032,han-lin-lecture-notes:page-0033}
\lean{HanLinLectureNotes.Ch02.nonpositive_subsolution_strong_alternative}
\leanok
If $u\in C^2(\Omega)\cap C(\overline\Omega)$ satisfies $Lu\ge0$ and $u\le0$ in
$\Omega$, then either $u<0$ throughout $\Omega$ or $u\equiv0$.
\end{theorem}
\begin{proof}
\sketch If $u(x_0)=0$ at an interior point, write $c=c^+-c^-$. Then
\[
a_{ij}D_{ij}u+b_iD_iu-c^-u\ge -c^+u\ge 0.
\]
Theorem~\ref{hanlin:maxprin-theorem-2-8} gives $u\equiv0$. Alternatively, for
$v=ue^{-\alpha x_1}$,
\[
a_{ij}D_{ij}v+\bigl(\alpha(a_{1i}+a_{i1})+b_i\bigr)D_iv+(a_{11}\alpha^2+b_1\alpha+c)v\ge 0.
\]
For large $\alpha$, the zero-order coefficient is positive; since $v\le0$, this
implies
\[
a_{ij}D_{ij}v+\bigl(\alpha(a_{1i}+a_{i1})+b_i\bigr)D_iv\ge 0.
\]
Theorem~\ref{hanlin:maxprin-theorem-2-8} applied to $v$ yields the dichotomy.
\end{proof}

\begin{theorem}\label{hanlin:maxprin-theorem-2-13}\dcref{Theorem 2.13}\group{ch02-maximum-principles}
\source{han-lin-lecture-notes:page-0033}
Assume that $w\in C^2(\Omega)\cap C^1(\overline\Omega)$ is positive on
$\overline\Omega$ and satisfies $Lw\le0$. If
$u\in C^2(\Omega)\cap C(\overline\Omega)$ and $Lu\ge0$, then a nonnegative
interior maximum of $u/w$ forces $u/w$ to be constant. If a nonconstant $u/w$
attains such a maximum at $x_0\in\partial\Omega$, then every outward direction
$\nu$ at $x_0$ satisfies
\[
\frac{\partial}{\partial\nu}\left(\frac{u}{w}\right)(x_0)>0,
\]
provided the interior sphere property holds at $x_0$.
\end{theorem}
\begin{proof}
\sketch For $v=u/w$,
\[
a_{ij}D_{ij}v+B_iD_iv+\left(\frac{Lw}{w}\right)v\ge 0,
\]
where $B_i=b_i+\dfrac{2}{w}a_{ij}D_jw$. Since $Lw/w\le0$, apply
Theorem~\ref{hanlin:maxprin-theorem-2-6} and
Corollary~\ref{hanlin:maxprin-cor-2-9} to $v$.
\end{proof}

\begin{remark}\label{hanlin:maxprin-remark-2-14}\dcref{Remark 2.14}\group{ch02-maximum-principles}
\source{han-lin-lecture-notes:page-0033}
\lean{HanLinLectureNotes.Ch02.positive_supersolution_comparison_and_dirichlet_uniqueness}
\leanok
Under the positive-supersolution hypothesis of
Theorem~\ref{hanlin:maxprin-theorem-2-13}, $L$ has the comparison principle.
Consequently,
\[
Lu=f\quad\text{in }\Omega,
\qquad
u=\varphi\quad\text{on }\partial\Omega
\]
has at most one solution.
\end{remark}

\begin{proposition}[Maximum principle on narrow domains]\label{hanlin:maxprin-prop-2-15}\dcref{Proposition 2.15}\group{ch02-maximum-principles}
\source{han-lin-lecture-notes:page-0033,han-lin-lecture-notes:page-0034}
\lean{HanLinLectureNotes.Ch02.exists_positive_supersolution_of_narrow_domain}
\leanok
Suppose a unit vector $e$ and $d>0$ satisfy
$|(y-x)\cdot e|<d$ for all $x,y\in\Omega$. There is $d_0>0$, depending only
on $\lambda$, $\|b_i\|_\infty$, and $\|c^+\|_\infty$, such that $d\le d_0$
implies the positive-supersolution hypothesis of
Theorem~\ref{hanlin:maxprin-theorem-2-13}.
\end{proposition}
\begin{proof}
\sketch Rotate and translate so that $\overline\Omega\subset\{0<x_1<d\}$,
and choose $N$ with $|b_i|,c^+\le N$. For
$w=e^{\alpha d}-e^{\alpha x_1}>0$,
\[
Lw=-(a_{11}\alpha^2+b_1\alpha)e^{\alpha x_1}+c(e^{\alpha d}-e^{\alpha x_1})
\le -(a_{11}\alpha^2+b_1\alpha)+Ne^{\alpha d}.
\]
Fix $\alpha$ so large that
\[
a_{11}\alpha^2+b_1\alpha\ge \lambda\alpha^2-N\alpha\ge 2N.
\]
Then $e^{\alpha d}\le2$ for sufficiently small $d$, and hence
$Lw\le N(e^{\alpha d}-2)\le0$.
\end{proof}

\begin{remark}\label{hanlin:maxprin-remark-2-16}\dcref{Remark 2.16}\group{ch02-maximum-principles}
\source{han-lin-lecture-notes:page-0034}
\lean{HanLinLectureNotes.Ch02.exists_positive_supersolution_of_narrow_unbounded_domain}
\leanok
Proposition~\ref{hanlin:maxprin-prop-2-15}, and some related results above, also
apply to unbounded domains.
\end{remark}

\section{Poisson Equations}\label{hanlin:maxprin-section-2-2}

\begin{lemma}\label{hanlin:poisson-lemma-2-17}\dcref{Lemma 2.17}\group{ch02-maximum-principles}
\source{han-lin-lecture-notes:page-0034,han-lin-lecture-notes:page-0035}
\lean{HanLinLectureNotes.Ch02.poisson_fderiv_norm_le}
\leanok
If $u\in C^2(B_R)\cap C(\overline B_R)$ satisfies $\Delta u=f$ in $B_R$ for
$f\in C(\overline B_R)$, then
\[
|Du(0)|\le \frac{n}{R}\max_{\partial B_R}|u|
 +\frac{R}{2}\max_{B_R}|f|.
\]
\end{lemma}
\begin{proof}
\sketch Put $M=\max_{\partial B_R}|u|$, $F=\max_{B_R}|f|$, and, on the upper
half-ball,
\[
v(x',x_n)=\frac{u(x',x_n)-u(x',-x_n)}2.
\]
Thus $|\Delta v|\le F$, $|v|\le M$ on the spherical boundary, and $v=0$ on
$x_n=0$. For
\[
w(x',x_n)=A|x'|^2+Bx_n+Cx_n^2
\]
choose
\[
A=\frac{M}{R^2},\qquad
C=-\frac F2-(n-1)A,\qquad
B=\frac R2F+\frac nR M.
\]
Then $\Delta w\le\Delta v$ and $w\ge v$ on the boundary. By
Corollary~\ref{hanlin:maxprin-cor-2-9}, $v\le w$. Setting $x'=0$, dividing by
$x_n$, and sending $x_n\downarrow0$ gives
\[
D_nu(0)\le\frac nR M+\frac R2F.
\]
Apply the same argument to $-u$ and rotate coordinates.
\end{proof}

\begin{lemma}\label{hanlin:poisson-lemma-2-18}\dcref{Lemma 2.18}\group{ch02-maximum-principles}
\source{han-lin-lecture-notes:page-0035,han-lin-lecture-notes:page-0036}
Let $u\in C^2(B_{2R})\cap C(\overline B_{2R})$ satisfy $\Delta u=f$ in
$B_{2R}$, where $f\in C(\overline B_{2R})$. For distinct $x,y\in B_R$,
\[
R^2\frac{|Du(x)-Du(y)|}{|x-y|}
\le c\left(\sup_{B_{2R}}|u|+R^2\sup_{B_{2R}}|f|\right)
\left(\log\frac{2R}{|x-y|}+1\right),
\]
where $c$ depends only on $n$.
\end{lemma}
\begin{proof}
\sketch It suffices to work with cubes $Q_R\subset Q_{2R}$. Write
\[
M=\sup_{Q_{2R}}|u|,\qquad F=\sup_{Q_{2R}}|f|.
\]
Lemma~\ref{hanlin:poisson-lemma-2-17} gives
\[
\sup_{Q_R}|Du|\le c\left(\frac MR+RF\right)=:\mu.
\]
On
\[
Q'=\left\{(x',y,z):|x_i|<\frac R2\ (1\le i<n),\quad
0<y,z<\frac R4\right\}\subset\mathbb R^{n+1},
\]
set
\[
v(x',y,z)=\frac14\{u(x',y+z)-u(x',y-z)-u(x',-y+z)+u(x',-y-z)\}
\]
and
\[
\mathcal L=\sum_{i=1}^{n-1}D_{x_ix_i}+\frac12D_{yy}+\frac12D_{zz}.
\]
Then $|\mathcal Lv|\le F$, $v=0$ when $y=0$ or $z=0$, $|v|\le M$ on
the lateral faces, and the remaining faces are bounded by $\mu z$ or $\mu y$.
For
\[
w(x',y,z)=\frac{4M|x'|^2}{R^2}
+yz\left(\frac{4\mu}{R}+k\log\frac{2R}{y+z}\right),
\qquad
k\ge\frac43\left(F+\frac{8(n-1)M}{R^2}\right),
\]
one has $\mathcal Lw\le-F$ and $w\ge|v|$ on $\partial Q'$. Hence
Corollary~\ref{hanlin:maxprin-cor-2-9} gives $|v|\le w$, and differentiation
at $z=0$ yields
\[
\frac12|D_yu(0,y)-D_yu(0,-y)|
\le\frac{4\mu}{R}y+ky\log\frac{2R}{y}.
\]
For the remaining derivative components, use
\[
Q'=\left\{(\widehat x,y,z):|x_i|<\frac R2\ (1\le i<n-1),\quad
0<y,z<\frac R2\right\}\subset\mathbb R^n
\]
and
\[
v(\widehat x,y,z)=\frac14\{u(\widehat x,y,z)-u(\widehat x,-y,z)
-u(\widehat x,y,-z)+u(\widehat x,-y,-z)\}.
\]
The comparison function
\[
w(\widehat x,y,z)=\frac{4M|\widehat x|^2}{R^2}
+yz\left(\frac{4\mu}{R}+\overline k\log\frac{2R}{y+z}\right),
\qquad
\overline k\ge\frac23\left(F+\frac{8(n-2)M}{R^2}\right),
\]
gives
\[
\frac12|D_{n-1}u(0,z)-D_{n-1}u(0,-z)|
\le\frac{4\mu}{R}z+\overline k z\log\frac{2R}{z},
\]
and likewise for $D_i$, $1\le i<n-1$. Translation, rotation, and scaling
give the stated estimate.
\end{proof}

The logarithmic modulus in Lemma~\ref{hanlin:poisson-lemma-2-18} cannot in
general be improved without stronger continuity of $f$.

For $0<\alpha<1$, use the pointwise seminorm
\[
[u]_{C^\alpha}(0)=\sup_{0<|x|\le1}\frac{|u(x)-u(0)|}{|x|^\alpha}.
\]
Thus $u$ is $C^\alpha$ at $0$ when $|u(x)-u(0)|\le c|x|^\alpha$, and it is
$C^{2,\alpha}$ there when some polynomial $P$ of degree at most two satisfies
\[
|u(x)-P(x)|\le c|x|^{2+\alpha}.
\]

\begin{lemma}\label{hanlin:maxprin-lemma-2-19}\dcref{Lemma 2.19}\group{ch02-maximum-principles}
\source{han-lin-lecture-notes:page-0037,han-lin-lecture-notes:page-0038}
Fix $\alpha\in(0,1)$. There are $c_0>0$, $\mu\in(0,1)$, and $\epsilon_0>0$,
depending only on $n$ and $\alpha$, with the following property. If
$u\in C^2(B_1)$, $\Delta u=f$ in $B_1$, $f\in C(B_1)$, and
$|u|\le1$, $|f|\le\epsilon_0$, then a harmonic quadratic polynomial
\[
p(x)=\frac12x^TAx+B\cdot x+C
\]
satisfies
\[
|u-p|\le\mu^{2+\alpha}\quad\text{on }B_\mu,
\qquad
|A|+|B|+|C|\le c_0.
\]
\end{lemma}
\begin{proof}
\sketch Let $v$ be harmonic in $B_1$ with $v=u$ on $\partial B_1$. Comparison
with the quadratic barrier gives
\[
|u(x)-v(x)|\le\frac{1-|x|^2}{2n}\sup_{B_1}|f|.
\]
The second-order Taylor polynomial $p$ of $v$ at $0$ is harmonic and has
universally bounded coefficients, while the interior derivative estimates give
\[
|u(x)-p(x)|
\le C|x|^3\sup_{B_{1/2}}|D^3v|+\frac1{2n}\sup_{B_1}|f|.
\]
Choose $\mu$ so the first term is at most $\frac12\mu^{2+\alpha}$ on $B_\mu$,
then choose $\epsilon_0$ so the second has the same bound.
\end{proof}

\begin{theorem}\label{hanlin:maxprin-theorem-2-20}\dcref{Theorem 2.20}\group{ch02-maximum-principles}
\source{han-lin-lecture-notes:page-0038,han-lin-lecture-notes:page-0039}
Let $u\in C^2(B_1)$ satisfy $\Delta u=f$ in $B_1$, where $f$ is
$C^\alpha$ at $0$. There is a quadratic polynomial
\[
p(x)=\frac12x^TAx+B\cdot x+C
\]
such that
\[
|u(x)-p(x)|\le c_0|x|^{2+\alpha}
\bigl(\|u\|_{L^\infty(B_1)}+|f(0)|+[f]_{C^\alpha}(0)\bigr)
\quad(x\in B_1),
\]
and
\[
|A|+|B|+|C|
\le c_0\bigl(\|u\|_{L^\infty(B_1)}+|f(0)|+[f]_{C^\alpha}(0)\bigr),
\]
where $c_0$ depends only on $n$ and $\alpha$.
\end{theorem}
\begin{proof}\uses{hanlin:maxprin-lemma-2-19}
\sketch Subtract $f(0)|x|^2/(2n)$ and normalize so that $f(0)=0$,
$\|u\|_\infty\le1$, and $[f]_{C^\alpha}(0)\le\epsilon_0$. Iterating
Lemma~\ref{hanlin:maxprin-lemma-2-19} produces harmonic polynomials
\[
P_k(x)=\frac12x^TA_kx+B_k\cdot x+C_k
\]
with
\[
\|u-P_k\|_{L^\infty(B_{\mu^k})}\le\mu^{(2+\alpha)k},
\]
and
\[
|A_{k+1}-A_k|\le c\mu^{\alpha k},\quad
|B_{k+1}-B_k|\le c\mu^{(1+\alpha)k},\quad
|C_{k+1}-C_k|\le c\mu^{(2+\alpha)k}.
\]
Indeed, for
\[
w(y)=\frac{(u-P_k)(\mu^ky)}{\mu^{(2+\alpha)k}},
\qquad
\Delta w(y)=\frac{f(\mu^ky)}{\mu^{\alpha k}},
\]
the lemma supplies a bounded harmonic quadratic $p_0$ with
$\|w-p_0\|_{L^\infty(B_\mu)}\le\mu^{2+\alpha}$, and one sets
\[
P_{k+1}(x)=P_k(x)+\mu^{(2+\alpha)k}p_0(x/\mu^k).
\]
The coefficients converge to those of $p$, and for $|x|\le\mu^k$,
\[
|P_k(x)-p(x)|
\le c\bigl(|x|^2\mu^{\alpha k}+|x|\mu^{(1+\alpha)k}
+\mu^{(2+\alpha)k}\bigr)
\le c\mu^{(2+\alpha)k}.
\]
Combining this with the approximation by $P_k$ and undoing the normalization
proves the estimates.
\end{proof}

\begin{theorem}\label{hanlin:maxprin-theorem-2-21}\dcref{Theorem 2.21}\group{ch02-maximum-principles}
\source{han-lin-lecture-notes:page-0040,han-lin-lecture-notes:page-0041}
Suppose $u\in C^2(\overline B_1)$ is positive and
\[
\Delta u=f(u)\quad\text{in }B_1,
\qquad
u=0\quad\text{on }\partial B_1,
\]
where $f$ is locally Lipschitz. Then $u$ is radial and
$\partial u/\partial r<0$ on $B_1\setminus\{0\}$.
\end{theorem}
\begin{proof}\uses{hanlin:maxprin-lemma-2-22,hanlin:maxprin-theorem-2-13,hanlin:maxprin-theorem-2-6}
\sketch In each direction, compare $u$ with its reflection across a moving
hyperplane. Lemma~\ref{hanlin:maxprin-lemma-2-22} starts the motion. At a first
stopping position, Theorem~\ref{hanlin:maxprin-theorem-2-13} and the Hopf lemma,
Theorem~\ref{hanlin:maxprin-theorem-2-6}, exclude contact. Reflection from both
directions gives radial symmetry, and the Hopf inequality gives strict radial
decrease.
\end{proof}

\begin{lemma}\label{hanlin:maxprin-lemma-2-22}\dcref{Lemma 2.22}\group{ch02-maximum-principles}
\source{han-lin-lecture-notes:page-0040}
\uses{hanlin:maxprin-theorem-2-21}
Let $u$ satisfy the hypotheses of Theorem~\ref{hanlin:maxprin-theorem-2-21}.
For $x_0\in\partial B_1$ and a unit vector $\nu$ with $\nu\cdot x_0>0$, there
is $r>0$ such that
\[
\frac{\partial u}{\partial\nu}(x)<0
\qquad(x\in B_1\cap B_r(x_0)).
\]
\end{lemma}
\begin{proof}
  \uses{hanlin:maxprin-theorem-2-6, hanlin:maxprin-lemma-2-22, hanlin:maxprin-theorem-2-13}
\sketch If $f(0)\le0$, local Lipschitz continuity gives a bounded $c\ge0$ with
\[
\Delta u=f(u)\le f(u)-f(0)\le c(x)u,
\]
so Theorem~\ref{hanlin:maxprin-theorem-2-6} gives the boundary derivative sign,
which persists nearby. If $f(0)>0$ and the conclusion fails along a sequence
approaching $x_0$, then $\partial_ru(x_0)=0$. Tangential differentiation of the
constant boundary value gives
\[
\frac{\partial^2u}{\partial T\partial r}(x_0)=0,
\qquad
\frac{\partial^2u}{\partial T^2}(x_0)=0,
\]
and hence
\[
\frac{\partial^2u}{\partial n^2}(x_0)=\Delta u(x_0)=f(0)>0.
\]
After rotating so $x_0=(1,0,\ldots,0)$, Taylor expansion gives
\[
\frac{\partial u}{\partial n}(x)
=\frac{\partial^2u}{\partial n^2}(x_0)(x_1-1)+o(|x-x_0|)<0
\]
for nearby interior points, a contradiction.
\end{proof}

\begin{remark}\label{hanlin:maxprin-remark-2-23}\dcref{Remark 2.23}\group{ch02-maximum-principles}
\source{han-lin-lecture-notes:page-0041}
The proof above uses boundary smoothness of both the domain and the solution.
Section~\ref{hanlin:maxprin-section-2-6} removes those requirements.
\end{remark}

\section{A Priori Estimates}\label{hanlin:maxprin-section-2-3}

Here $\Omega\subset\mathbb R^n$ is bounded and connected,
\[
Lu=a_{ij}(x)D_{ij}u+b_i(x)D_iu+c(x)u,
\]
the coefficients are continuous on $\overline\Omega$, and
\[
a_{ij}(x)\xi_i\xi_j\ge\lambda|\xi|^2,
\qquad
\max_\Omega|a_{ij}|+\max_\Omega|b_i|\le\Lambda.
\]

\begin{theorem}\label{hanlin:maxprin-theorem-2-24}\dcref{Theorem 2.24}\group{ch02-maximum-principles}
\source{han-lin-lecture-notes:page-0042}
\lean{HanLinLectureNotes.Ch02.dirichlet_apriori_estimate}
\leanok
Let $u\in C^2(\Omega)\cap C(\overline\Omega)$ solve
\[
Lu=f\quad\text{in }\Omega,
\qquad
u=\varphi\quad\text{on }\partial\Omega,
\]
where $f\in C(\overline\Omega)$, $\varphi\in C(\partial\Omega)$, and $c\le0$.
Then
\[
\max_\Omega|u|\le\max_{\partial\Omega}|\varphi|+C\max_\Omega|f|,
\]
where $C$ depends only on $\lambda$, $\Lambda$, and
$\operatorname{diam}(\Omega)$.
\end{theorem}
\begin{proof}\uses{hanlin:maxprin-cor-2-9}
\sketch Put $F=\max_\Omega|f|$ and
$\Phi=\max_{\partial\Omega}|\varphi|$. After a rigid motion, suppose
$\Omega\subset\{0<x_1<d\}$, and set
\[
w=\Phi+(e^{\alpha d}-e^{\alpha x_1})F.
\]
Choose $\alpha$ so that $\lambda\alpha^2+b_1(x)\alpha\ge1$ in $\Omega$.
Since $c\le0$,
\[
-Lw=(a_{11}\alpha^2+b_1\alpha)Fe^{\alpha x_1}
-c\Phi-c(e^{\alpha d}-e^{\alpha x_1})F\ge F.
\]
Thus $L(w\pm u)\le0$ and $w\pm u\ge0$ on the boundary.
Corollary~\ref{hanlin:maxprin-cor-2-9} yields
\[
|u|\le w\le\Phi+(e^{\alpha d}-1)F.
\]
\end{proof}

\begin{theorem}\label{hanlin:maxprin-theorem-2-25}\dcref{Theorem 2.25}\group{ch02-maximum-principles}
\source{han-lin-lecture-notes:page-0043,han-lin-lecture-notes:page-0044}
Let $u\in C^2(\Omega)\cap C^1(\overline\Omega)$ solve
\[
Lu=f\quad\text{in }\Omega,
\qquad
\frac{\partial u}{\partial n}+\alpha(x)u=\varphi
\quad\text{on }\partial\Omega.
\]
If $c\le0$ in $\Omega$ and $\alpha\ge\alpha_0>0$ on $\partial\Omega$, then
\[
\max_{\partial\Omega}|u|
\le C\left(\max_{\partial\Omega}|\varphi|+\max_\Omega|f|\right),
\]
where $C$ depends only on $\lambda$, $\Lambda$, $\alpha_0$, and
$\operatorname{diam}(\Omega)$.
\end{theorem}
\begin{proof}
\sketch First assume $c\le-c_0<0$. With
$F=\max_\Omega|f|$, $\Phi=\max_{\partial\Omega}|\varphi|$, set
\[
v=\frac{F}{c_0}+\frac{\Phi}{\alpha_0}\pm u.
\]
Then
\[
Lv\le-F\pm f\le0,
\qquad
\frac{\partial v}{\partial n}+\alpha v\ge\Phi\pm\varphi\ge0.
\]
If $v$ had a negative minimum, Lemma~\ref{hanlin:maxprin-lemma-2-3} would place
it at $x_0\in\partial\Omega$, where $\partial_nv(x_0)\le0$ and therefore
$(\partial_nv+\alpha v)(x_0)<0$, a contradiction. Hence
\[
|u|\le\frac{F}{c_0}+\frac{\Phi}{\alpha_0}.
\]

For general $c\le0$, write $u=zw$ with $z>0$. Then
\[
a_{ij}D_{ij}w+B_iD_iw+
\left(c+\frac{a_{ij}D_{ij}z+b_iD_iz}{z}\right)w=\frac fz,
\]
\[
\frac{\partial w}{\partial n}
+\left(\alpha+\frac1z\frac{\partial z}{\partial n}\right)w
=\frac\varphi z,
\qquad
B_i=\frac1z(a_{ij}+a_{ji})D_jz+b_i.
\]
If $\Omega\subset\{0<x_1<d\}$, choose
\[
z(x)=A+e^{\beta d}-e^{\beta x_1}.
\]
Taking $\beta$ so that $\beta^2a_{11}+\beta b_1\ge1$ gives
\[
-\frac{a_{ij}D_{ij}z+b_iD_iz}{z}
=\frac{(\beta^2a_{11}+\beta b_1)e^{\beta x_1}}
{A+e^{\beta d}-e^{\beta x_1}}
\ge\frac1{A+e^{\beta d}}>0.
\]
Taking $A$ large also ensures
\[
\left|\frac1z\frac{\partial z}{\partial n}\right|
\le\frac\beta A e^{\beta d}\le\frac{\alpha_0}{2}.
\]
The equation for $w$ is therefore in the strict case above, with controlled
coefficients; boundedness above and below of $z$ gives the result for $u$.
\end{proof}

\begin{remark}\label{hanlin:maxprin-remark-2-26}\dcref{Remark 2.26}\group{ch02-maximum-principles}
\source{han-lin-lecture-notes:page-0044}
\lean{HanLinLectureNotes.zero_robin_coefficient_nonunique}
\leanok
The conclusion may fail under the weaker assumption $\alpha\ge0$; even
uniqueness is then not guaranteed.
\end{remark}

\section{Gradient Estimates}\label{hanlin:maxprin-section-2-4}

Consider
\[
a_{ij}(x)D_{ij}u+b_i(x)D_iu=f(x,u)
\quad\text{in }\Omega,
\]
where $\Omega\subset\mathbb R^n$ is bounded and connected, the coefficients are
bounded and continuous, and
\[
a_{ij}(x)\xi_i\xi_j\ge\lambda|\xi|^2.
\]

\begin{theorem}\label{hanlin:maxprin-theorem-2-27}\dcref{Theorem 2.27}\group{ch02-maximum-principles}
\source{han-lin-lecture-notes:page-0044,han-lin-lecture-notes:page-0045}
Let $u\in C^3(\Omega)\cap C^1(\overline\Omega)$ satisfy the equation above, with
$a_{ij},b_i\in C^1(\overline\Omega)$ and
$f\in C^1(\overline\Omega\times\mathbb R)$. Then
\[
\sup_\Omega|Du|\le\sup_{\partial\Omega}|Du|+C,
\]
where, for $M=\|u\|_{L^\infty(\Omega)}$, the constant $C$ depends only on
$\lambda$, $\operatorname{diam}(\Omega)$, the $C^1$ norms of $a_{ij},b_i$, $M$,
and $\|f\|_{C^1(\overline\Omega\times[-M,M])}$.
\end{theorem}
\begin{proof}
\sketch Write $L=a_{ij}D_{ij}+b_iD_i$. Differentiating the equation and using
\[
D_i(|Du|^2)=2D_kuD_{ki}u,
\qquad
D_{ij}(|Du|^2)=2D_{ki}uD_{kj}u+2D_kuD_{kij}u,
\]
gives
\[
\begin{split}
L(|Du|^2)={}&2a_{ij}D_{ki}uD_{kj}u
-2D_ka_{ij}D_kuD_{ij}u-2D_kb_iD_kuD_iu\\
&+2D_zf|Du|^2+2D_kfD_ku.
\end{split}
\]
Ellipticity and Cauchy's inequality imply
\[
L(|Du|^2)\ge\lambda|D^2u|^2-C|Du|^2-C,
\]
whereas
\[
L(u^2)=2a_{ij}D_iuD_ju+2uf
\ge2\lambda|Du|^2+2uf.
\]
Choose $\alpha$ large, translate so $x_1>0$ on $\Omega$, and then choose
$\beta$ large. The function
\[
w=|Du|^2+\alpha u^2+e^{\beta x_1}
\]
satisfies
\[
Lw\ge\lambda|D^2u|^2+|Du|^2
+e^{\beta x_1}(\beta^2a_{11}+\beta b_1)-C\ge0.
\]
The weak maximum principle bounds $w$ by its boundary values and gives the
claimed estimate.
\end{proof}

\begin{proposition}\label{hanlin:maxprin-prop-2-28}\dcref{Proposition 2.28}\group{ch02-maximum-principles}
\source{han-lin-lecture-notes:page-0046}
Under the equation and coefficient hypotheses of
Theorem~\ref{hanlin:maxprin-theorem-2-27}, but with only
$u\in C^3(\Omega)$, every $\Omega'\Subset\Omega$ satisfies
\[
\sup_{\Omega'}|Du|\le C.
\]
Here $C$ depends on $\lambda$, $\operatorname{diam}(\Omega)$,
$\operatorname{dist}(\Omega',\partial\Omega)$, the indicated coefficient
$C^1$ norms, $M=\|u\|_{L^\infty(\Omega)}$, and
$\|f\|_{C^1(\Omega\times[-M,M])}$.
\end{proposition}
\begin{proof}
\sketch Choose $\gamma\in C_c^\infty(\Omega)$, $\gamma\ge0$, with
$|D\gamma|^2\le C\gamma$, and set
\[
v=\gamma|Du|^2,
\qquad
w=v+\alpha u^2+e^{\beta x_1}.
\]
The calculation from Theorem~\ref{hanlin:maxprin-theorem-2-27} gives
\[
Lv=(L\gamma)|Du|^2+\gamma L(|Du|^2)
+2a_{ij}D_i\gamma D_j(|Du|^2).
\]
Using
\[
|4a_{ij}D_kuD_i\gamma D_{kj}u|
\le\varepsilon|D\gamma|^2|D^2u|^2+C(\varepsilon)|Du|^2
\]
and small $\varepsilon$ yields
\[
Lv\ge\frac\lambda2\gamma|D^2u|^2-C|Du|^2-C.
\]
Large $\alpha,\beta$ make $Lw\ge0$. Since $\gamma$ vanishes near the boundary,
the maximum principle bounds $\gamma|Du|^2$, and a cutoff equal to one on
$\Omega'$ proves the estimate.
\end{proof}

\begin{proposition}\label{hanlin:maxprin-prop-2-29}\dcref{Proposition 2.29}\group{ch02-maximum-principles}
\source{han-lin-lecture-notes:page-0047}
Assume that $\Omega$ has the uniform exterior sphere property and that
$u\in C^2(\Omega)\cap C(\overline\Omega)$ satisfies
\[
a_{ij}(x)D_{ij}u+b_i(x)D_iu=f(x,u)
\quad\text{in }\Omega,
\]
with continuous $a_{ij},b_i$ and $f$. Then
\[
|u(x)-u(x_0)|\le C|x-x_0|
\qquad(x\in\Omega,\ x_0\in\partial\Omega).
\]
If $M=\|u\|_{L^\infty(\Omega)}$ and
$\varphi\in C^2(\overline\Omega)$ satisfies $\varphi=u$ on $\partial\Omega$,
then $C$ depends only on $\lambda$, $\Omega$, the coefficient sup norms, $M$,
$\|f\|_{L^\infty(\Omega\times[-M,M])}$, and
$\|\varphi\|_{C^2(\overline\Omega)}$.
\end{proposition}
\begin{proof}
\sketch Subtract the extension $\varphi$, so it is enough to treat $u=0$ on
$\partial\Omega$. Write $L=a_{ij}D_{ij}+b_iD_i$ and
$F=\sup_\Omega|f(\cdot,u)|$. Fix $x_0\in\partial\Omega$ and an exterior ball
$B_R(y)$ tangent at $x_0$. For
\[
d(x)=|x-y|-R,
\qquad 0<d(x)<D:=\operatorname{diam}(\Omega),
\]
seek a barrier $w=\psi(d)$ with $w(x_0)=0$, $w>0$ elsewhere, and $Lw\le-F$.
The distance derivatives satisfy
\[
D_id=\frac{x_i-y_i}{|x-y|},
\qquad
D_{ij}d=\frac{\delta_{ij}}{|x-y|}
-\frac{(x_i-y_i)(x_j-y_j)}{|x-y|^3},
\]
and, if $\Lambda=\sup|a_{ij}|$,
\[
|Dd|=1,
\qquad
a_{ij}D_{ij}d\le\frac{n\Lambda-\lambda}{R}.
\]
For $\psi'>0$ and $\psi''<0$ this gives
\[
Lw\le\lambda\psi''
+\left(\frac{n\Lambda-\lambda}{R}+|b|\right)\psi'.
\]
Choose positive constants $a,b$ so that any solution of
\[
\psi''+a\psi'+b=0
\]
satisfies the required inequality, and take
\[
\psi(d)=\frac ba\left\{\frac1a e^{aD}(1-e^{-ad})-d\right\}.
\]
Then
\[
\psi(0)=0,
\qquad
\psi'(d)=\frac ba(e^{a(D-d)}-1)>0\quad(0\le d<D),
\qquad
\psi''(d)=-be^{a(D-d)}<0.
\]
Thus $Lw\le-F$, and comparison gives $-w\le u\le w$. Since $\psi'$ is
uniformly bounded and $d(x)\le|x-x_0|$, the desired Lipschitz bound follows;
the initial subtraction restores the dependence on $\varphi$.
\end{proof}

\section{Alexandroff Maximum Principle}\label{hanlin:maxprin-section-2-5}

Let $\Omega\subset\mathbb R^n$ be bounded and
\[
L=a_{ij}(x)D_{ij}+b_i(x)D_i+c(x),
\]
where the coefficients are continuous and $A=(a_{ij})$ is positive definite.
Write
\[
D=\det A,
\qquad
D^*=D^{1/n},
\qquad
0<\lambda\le D^*\le\Lambda.
\]
For $u\in C(\overline\Omega)$, define its upper contact set by
\[
\Gamma^+=\left\{y\in\Omega:
u(x)\le u(y)+p\cdot(x-y)\ \text{for all }x\in\Omega
\text{ and some }p\in\mathbb R^n\right\}.
\]
If $u\in C^1$, then $p=Du(y)$; if $u\in C^2$, then $D^2u\le0$ on
$\Gamma^+$. Moreover, $u$ is concave exactly when $\Gamma^+=\Omega$.

\begin{theorem}[Alexandroff maximum principle]\label{hanlin:maxprin-theorem-2-30}\dcref{Theorem 2.30}\group{ch02-maximum-principles}
\source{han-lin-lecture-notes:page-0049,han-lin-lecture-notes:page-0050,han-lin-lecture-notes:page-0051}
Suppose $u\in C(\overline\Omega)\cap C^2(\Omega)$ satisfies $Lu\ge f$, $c\le0$,
and
\[
\frac{|b|}{D^*},\ \frac f{D^*}\in L^n(\Omega).
\]
Then
\[
\sup_\Omega u\le\sup_{\partial\Omega}u^+
+C\left\|\frac{f^-}{D^*}\right\|_{L^n(\Gamma^+)},
\]
where $C$ depends only on $n$, $d=\operatorname{diam}(\Omega)$, and
$\||b|/D^*\|_{L^n(\Gamma^+)}$. One admissible value is
\[
C=d\left\{\exp\!\left(
\frac{2^{n-2}}{\omega_n^n}
\left(\left\|\frac b{D^*}\right\|_{L^n(\Gamma^+)}+1\right)
\right)-1\right\},
\]
where $\omega_n$ is the volume of the unit ball in $\mathbb R^n$.
\end{theorem}
\begin{proof}
\uses{hanlin:maxprin-lemma-2-33}
\sketch Apply Lemma~\ref{hanlin:maxprin-lemma-2-33} with a regularized
inverse-power weight. On $\Gamma^+$, the equation and the determinant--trace
inequality bound the Jacobian of the gradient map by $f^-$ and the drift term.
Integration and passage to the zero-regularization limit give the displayed
constant and estimate.
\end{proof}

\begin{remark}\label{hanlin:maxprin-remark-2-31}\dcref{Remark 2.31}\group{ch02-maximum-principles}
\source{han-lin-lecture-notes:page-0049}
The $L^n$ norm of $f^-/D^*$ may be restricted to
\[
\Gamma^+\cap\left\{x\in\Omega:
u(x)>\sup_{\partial\Omega}u^+\right\}.
\]
\end{remark}

\begin{remark}\label{hanlin:maxprin-remark-2-32}\dcref{Remark 2.32}\group{ch02-maximum-principles}
\source{han-lin-lecture-notes:page-0049}
Unlike Theorem~\ref{hanlin:maxprin-theorem-2-24}, this result does not assume
uniform ellipticity.
\end{remark}

\begin{lemma}\label{hanlin:maxprin-lemma-2-33}\dcref{Lemma 2.33}\group{ch02-maximum-principles}
\source{han-lin-lecture-notes:page-0049,han-lin-lecture-notes:page-0050}
\lean{HanLinLectureNotes.Ch02.alexandroff_weighted_contact_integral}
\leanok
Let $g\in L^1_{\mathrm{loc}}(\mathbb R^n)$ be nonnegative and
$u\in C(\overline\Omega)\cap C^2(\Omega)$. If
\[
\widetilde M=\frac{\sup_\Omega u-\sup_{\partial\Omega}u^+}
{d},
\qquad d=\operatorname{diam}(\Omega),
\]
then
\[
\int_{B_{\widetilde M}(0)}g
\le\int_{\Gamma^+}g(Du)|\det D^2u|.
\]
\end{lemma}
\begin{proof}
\sketch Every $p\in B_{\widetilde M}(0)$ is the slope of an affine function
supporting $u$ from above at some point of $\Gamma^+$. Thus
$B_{\widetilde M}(0)\subset Du(\Gamma^+)$, and the area formula for the gradient
map gives the inequality.
\end{proof}

\begin{remark}\label{hanlin:maxprin-remark-2-34}\dcref{Remark 2.34}\group{ch02-maximum-principles}
\source{han-lin-lecture-notes:page-0050}
\uses{hanlin:maxprin-lemma-2-33}
\lean{HanLinLectureNotes.Ch02.alexandroff_operator_contact_inequalities}
\leanok
For positive definite $A=(a_{ij})$, the arithmetic--geometric mean inequality
on $\GammaP$ gives
\[
\det(-D^2u)\le\frac1D
\left(\frac{-a_{ij}D_{ij}u}{n}\right)^n.
\]
Consequently, Lemma~\ref{hanlin:maxprin-lemma-2-33} implies
\[
\int_{B_{\Mtilde}(0)}g
\le\int_{\GammaP}g(Du)
\left(\frac{-a_{ij}D_{ij}u}{nD^*}\right)^n.
\]
\end{remark}

\begin{remark}\label{hanlin:maxprin-remark-2-35}\dcref{Remark 2.35}\group{ch02-maximum-principles}
\source{han-lin-lecture-notes:page-0050}
\uses{hanlin:maxprin-lemma-2-33}
\lean{HanLinLectureNotes.Ch02.alexandroff_zero_drift_contact_bounds}
\leanok
Taking $g=1$ in Lemma~\ref{hanlin:maxprin-lemma-2-33} yields
\[
\sup_\Omega u\le\sup_{\partial\Omega}u^+
+\frac d{\omega_n^{1/n}}
\left(\int_{\GammaP}|\det D^2u|\right)^{1/n}
\]
\[
\le\sup_{\partial\Omega}u^+
+\frac d{\omega_n^{1/n}}
\left(\int_{\GammaP}
\left(\frac{-a_{ij}D_{ij}u}{nD^*}\right)^n\right)^{1/n}.
\]
This is the zero-drift, zero-order-free case of
Theorem~\ref{hanlin:maxprin-theorem-2-30}.
\end{remark}

\begin{theorem}\label{hanlin:maxprin-theorem-2-36}\dcref{Theorem 2.36}\group{ch02-maximum-principles}
\source{han-lin-lecture-notes:page-0051,han-lin-lecture-notes:page-0052}
Suppose $u\in C(\overline\Omega)\cap C^2(\Omega)$ solves
\[
Qu=a_{ij}(x,u,Du)D_{ij}u+b(x,u,Du)=0
\quad\text{in }\Omega,
\]
where $a_{ij}$ is continuous on
$\Omega\times\mathbb R\times\mathbb R^n$ and
\[
a_{ij}(x,z,p)\xi_i\xi_j>0
\quad\text{for }(x,z,p)\in\Omega\times\mathbb R\times\mathbb R^n,
\ \xi\ne0.
\]
Assume nonnegative functions $g\in L^n_{\mathrm{loc}}(\mathbb R^n)$ and
$h\in L^n(\Omega)$ satisfy
\[
\frac{|b(x,z,p)|}{nD^*}\le\frac{h(x)}{g(p)}
\]
for all $(x,z,p)$ and
\[
\int_\Omega h^n(x)\,dx
<\int_{\mathbb R^n}g^n(p)\,dp=:g_\infty.
\]
Then
\[
\sup_\Omega|u|\le\sup_{\partial\Omega}|u|+C\operatorname{diam}(\Omega),
\]
where $C$ depends only on $g$ and $h$.
\end{theorem}
\begin{proof}
\uses{hanlin:maxprin-lemma-2-33}
\sketch Apply Lemma~\ref{hanlin:maxprin-lemma-2-33} separately to $u$ and
$-u$, using $g^n$ as the weight. The structural inequality bounds the
contact-set Jacobian integral by $\int_\Omega h^n$. Its strict separation from
$g_\infty$ bounds the attainable supporting slopes and hence the oscillation by
$C\operatorname{diam}(\Omega)$.
\end{proof}

\begin{example}\label{hanlin:maxprin-example-2-37}\dcref{Example 2.37}\group{ch02-maximum-principles}
\source{han-lin-lecture-notes:page-0052}
\lean{HanLinLectureNotes.Ch02.prescribedMeanCurvatureExample}
\leanok
For the prescribed mean-curvature equation
\[
(1+|Du|^2)\Delta u-D_iuD_juD_{ij}u
=nH(x)(1+|Du|^2)^{3/2},
\]
the coefficients in Theorem~\ref{hanlin:maxprin-theorem-2-36} are
\[
a_{ij}(x,z,p)=(1+|p|^2)\delta_{ij}-p_ip_j,
\qquad
b(x,z,p)=-nH(x)(1+|p|^2)^{3/2}.
\]
Thus $D=(1+|p|^2)^{n-1}$ and
\[
\frac{|b(x,z,p)|}{nD^*}
=|H(x)|(1+|p|^2)^{(n+2)/(2n)}.
\]
The corresponding weight satisfies
\[
g_\infty=\int_{\mathbb R^n}\frac{dp}{(1+|p|^2)^{(n+2)/2}}
=\omega_n.
\]
\end{example}

\begin{corollary}\label{hanlin:maxprin-cor-2-38}\dcref{Corollary 2.38}\group{ch02-maximum-principles}
\source{han-lin-lecture-notes:page-0052}
Let $u\in C(\overline\Omega)\cap C^2(\Omega)$ satisfy
\[
(1+|Du|^2)\Delta u-D_iuD_juD_{ij}u
=nH(x)(1+|Du|^2)^{3/2}
\quad\text{in }\Omega,
\]
where $H\in C(\Omega)$ and
\[
H_0=\int_\Omega|H(x)|^n\,dx<\omega_n.
\]
Then
\[
\sup_\Omega|u|\le\sup_{\partial\Omega}|u|
+C\operatorname{diam}(\Omega),
\]
where $C$ depends only on $n$ and $H_0$.
\end{corollary}

\begin{theorem}\label{hanlin:maxprin-theorem-2-39}\dcref{Theorem 2.39}\group{ch02-maximum-principles}
\source{han-lin-lecture-notes:page-0053}
Suppose $u\in C(\overline\Omega)\cap C^2(\Omega)$ satisfies
\[
\det(D^2u)=f(x,u,Du)
\quad\text{in }\Omega,
\]
where $f\in C(\Omega\times\mathbb R\times\mathbb R^n)$. Assume nonnegative
$g\in L^1_{\mathrm{loc}}(\mathbb R^n)$ and $h\in L^1(\Omega)$ satisfy
\[
|f(x,z,p)|\le\frac{h(x)}{g(p)}
\]
for all $(x,z,p)$ and
\[
\int_\Omega h(x)\,dx
<\int_{\mathbb R^n}g(p)\,dp=:g_\infty.
\]
Then
\[
\sup_\Omega|u|\le\sup_{\partial\Omega}|u|
+C\operatorname{diam}(\Omega),
\]
where $C$ depends only on $g$ and $h$.
\end{theorem}
\begin{proof}
\sketch Apply the contact-set estimate of
Lemma~\ref{hanlin:maxprin-lemma-2-33} to $u$ and $-u$. The equation and the
bound $|f|\le h/g$ control the weighted gradient-image measure by
$\int_\Omega h$, while the strict inequality with $g_\infty$ bounds the support
slopes. Multiplication by $\operatorname{diam}(\Omega)$ gives the estimate.
\end{proof}

\begin{corollary}\label{hanlin:maxprin-cor-2-40}\dcref{Corollary 2.40}\group{ch02-maximum-principles}
\source{han-lin-lecture-notes:page-0053}
If $u\in C(\overline\Omega)\cap C^2(\Omega)$ and
\[
\det(D^2u)=f(x)\quad\text{in }\Omega,
\qquad f\in C(\overline\Omega),
\]
then
\[
\sup_\Omega|u|\le\sup_{\partial\Omega}|u|
+\frac{\operatorname{diam}(\Omega)}{\omega_n^{1/n}}
\left(\int_\Omega|f|^n\right)^{1/n}.
\]
\end{corollary}

\begin{corollary}\label{hanlin:maxprin-cor-2-41}\dcref{Corollary 2.41}\group{ch02-maximum-principles}
\source{han-lin-lecture-notes:page-0053}
Let $u\in C(\overline\Omega)\cap C^2(\Omega)$ satisfy the prescribed
Gauss-curvature equation
\[
\det(D^2u)=K(x)(1+|Du|^2)^{(n+2)/2}
\quad\text{in }\Omega,
\]
where $K\in C(\overline\Omega)$ and
\[
K_0:=\int_\Omega|K(x)|\,dx<\omega_n.
\]
Then
\[
\sup_\Omega|u|\le\sup_{\partial\Omega}|u|
+C\operatorname{diam}(\Omega),
\]
where $C$ depends only on $n$ and $K_0$.
\end{corollary}

For the final result of this section, let
\[
Lu=a_{ij}D_{ij}u+b_iD_iu+cu,
\qquad
|b_i|+|c|\le\Lambda,
\qquad
\det(a_{ij})\ge\lambda>0,
\]
with $(a_{ij})$ pointwise positive definite.

\begin{theorem}\label{hanlin:maxprin-theorem-2-42}\dcref{Theorem 2.42}\group{ch02-maximum-principles}
\source{han-lin-lecture-notes:page-0054}
Suppose $u\in C(\overline\Omega)\cap C^2(\Omega)$ satisfies $Lu\ge0$ in
$\Omega$, $u\le0$ on $\partial\Omega$, and
$\operatorname{diam}(\Omega)\le d$. There is
$\delta>0$, depending only on $n$, $\lambda$, $\Lambda$, and $d$, such that
\[
|\Omega|\le\delta\quad\Longrightarrow\quad u\le0\text{ in }\Omega.
\]
\end{theorem}
\begin{proof}
\sketch The case $c\le0$ follows from
Theorem~\ref{hanlin:maxprin-theorem-2-30}. In general, write $c=c^+-c^-$, so
\[
a_{ij}D_{ij}u+b_iD_iu-c^-u\ge-c^+u.
\]
Applying Theorem~\ref{hanlin:maxprin-theorem-2-30} to the positive part gives
\[
\sup_\Omega u
\le C\|c^+u\|_{L^n(\Omega)}
\le C\|c^+\|_{L^\infty(\Omega)}|\Omega|^{1/n}\sup_\Omega u.
\]
Choose $\delta$ so the coefficient on the right is at most $1/2$.
\end{proof}

\section{The Moving Plane Method}\label{hanlin:maxprin-section-2-6}

\begin{theorem}\label{hanlin:maxprin-theorem-2-43}\dcref{Theorem 2.43}\group{ch02-maximum-principles}
\source{han-lin-lecture-notes:page-0054,han-lin-lecture-notes:page-0055}
Suppose $u\in C(\overline B_1)\cap C^2(B_1)$ is positive and satisfies
\[
\Delta u+f(u)=0\quad\text{in }B_1,
\qquad
u=0\quad\text{on }\partial B_1,
\]
where $f$ is locally Lipschitz on $\mathbb R$. Then $u$ is radial and
\[
\frac{\partial u}{\partial r}(x)<0
\qquad(x\in B_1\setminus\{0\}).
\]
\end{theorem}
\begin{proof}
\sketch Apply Lemma~\ref{hanlin:maxprin-lemma-2-44} to $B_1$ after every
rotation of coordinates. Symmetry across every hyperplane through the origin
implies radial symmetry, and the directional inequality gives strict radial
decrease away from the origin.
\end{proof}

\begin{lemma}\label{hanlin:maxprin-lemma-2-44}\dcref{Lemma 2.44}\group{ch02-maximum-principles}
\source{han-lin-lecture-notes:page-0055}
Let the bounded domain $\Omega$ be convex in the $x_1$ direction and symmetric
about $\{x_1=0\}$. If $u\in C(\overline\Omega)\cap C^2(\Omega)$ is positive and
\[
\Delta u+f(u)=0\quad\text{in }\Omega,
\qquad
u=0\quad\text{on }\partial\Omega,
\]
where $f$ is locally Lipschitz, then
\[
u(x_1,y)=u(-x_1,y),
\qquad
D_{x_1}u(x)<0\quad\text{when }x_1>0.
\]
\end{lemma}
\begin{proof}
\sketch Put
\[
a=\sup\{x_1:(x_1,y)\in\Omega\},
\qquad
\Sigma_\lambda=\{x\in\Omega:x_1>\lambda\},
\qquad
T_\lambda=\{x_1=\lambda\},
\]
and let $\Sigma'_\lambda$ be the reflection of $\Sigma_\lambda$ across
$T_\lambda$. For
\[
x_\lambda=(2\lambda-x_1,x_2,\ldots,x_n),
\qquad
w_\lambda(x)=u(x)-u(x_\lambda),
\]
local Lipschitz continuity of $f$ gives a bounded $c(x,\lambda)$ such that
\[
\Delta w_\lambda+c(x,\lambda)w_\lambda=0
\quad\text{in }\Sigma_\lambda,
\qquad
w_\lambda\le0\text{ on }\partial\Sigma_\lambda.
\]
For $\lambda$ close to $a$, either the narrow-domain principle,
Proposition~\ref{hanlin:maxprin-prop-2-15}, or the small-volume principle,
Theorem~\ref{hanlin:maxprin-theorem-2-42}, gives $w_\lambda<0$.

Let $(\lambda_0,a)$ be the maximal interval on which this strict inequality
holds. If $\lambda_0>0$, continuity and
Theorem~\ref{hanlin:maxprin-theorem-2-8} give
$w_{\lambda_0}<0$ in $\Sigma_{\lambda_0}$. Choose compact
$K\subset\Sigma_{\lambda_0}$ with
\[
|\Sigma_{\lambda_0}\setminus K|<\frac\delta2,
\qquad
w_{\lambda_0}\le-\eta<0\quad\text{on }K,
\]
where $\delta$ is supplied by Theorem~\ref{hanlin:maxprin-theorem-2-42}.
For small $\varepsilon>0$, continuity gives
$w_{\lambda_0-\varepsilon}<0$ on $K$, while
\[
|\Sigma_{\lambda_0-\varepsilon}\setminus K|<\delta.
\]
Theorem~\ref{hanlin:maxprin-theorem-2-42} yields
$w_{\lambda_0-\varepsilon}\le0$ on the complement, and
Theorem~\ref{hanlin:maxprin-theorem-2-12} makes the inequality strict. This
contradicts maximality, so $\lambda_0=0$.

Hence, whenever $x_1>0$, $x_1^*<x_1$, and $x_1+x_1^*>0$,
\[
u(x_1,y)<u(x_1^*,y).
\]
Letting $x_1^*\downarrow-x_1$ and repeating from the opposite direction gives
symmetry. Finally, Theorem~\ref{hanlin:maxprin-theorem-2-6} on $T_\lambda$ gives
\[
D_{x_1}w_\lambda|_{x_1=\lambda}
=2D_{x_1}u|_{x_1=\lambda}<0,
\]
which proves strict monotonicity.
\end{proof}
